\documentclass[12pt,a4paper]{report}
\usepackage[utf8]{inputenc}
\usepackage[T1]{fontenc}
\usepackage[spanish]{babel}
\usepackage{graphicx}
\usepackage{geometry}
\usepackage{tikz}
\usepackage{tikz}
\usetikzlibrary{calc}

% Márgenes
\geometry{
  top=2.5cm,
  bottom=2.5cm,
  left=3cm,
  right=2.5cm
}

%%%%%%%%%%%%%%%%%%%%%%%%%%%%%%%%%%%%%%%%%%%%%%%%%%%%%%%%%
%%%%%%%%%%%%% PIE DE PAGINA %%%%%%%%%%%%%%%%%%%%%%%%%%%%%
%%%%%%%%%%%%%%%%%%%%%%%%%%%%%%%%%%%%%%%%%%%%%%%%%%%%%%%%%
% Definir el estilo de página personalizado

\usepackage{fancyhdr}
\pagestyle{fancy}
\fancypagestyle{mypagestyle}{
    \fancyfoot{}                       % Limpia cualquier configuración previa
    \fancyfoot[R]{\thepage}       % Número de página en el lado derecho
}

% Aplicar el estilo de página personalizado a todas las páginas
\pagestyle{mypagestyle}
\setlength{\footskip}{-3cm}

%%%%%%%%%%%%%%%%%%%%%%%%%%%%%%%%%%%%%%%%%%%%%%%%%%%%%%%%%%
%%% ENCABEZADO DE LAS PÁGINAS TIPO UNICAFAM %%%%%%%%%%%%%%
%%%%%%%%%%%%%%%%%%%%%%%%%%%%%%%%%%%%%%%%%%%%%%%%%%%%%%%%%%
\setlength{\hoffset}{-0.5cm} 
% Ajustar margen superior
\setlength{\topmargin}{-1cm}  % mueve el encabezado más arriba
\newcommand{\MYhead}{\smash{\scriptsize
\hfil\parbox[t][\height][t]{\textwidth}{\centering
\begin{picture}(0,0) 
\put(-150,-19){\includegraphics[width=1.5cm, height=1.5cm]{logos/logo_ipn.png}} \end{picture} 
Prototipo de un sistema telem\'atico \\
\hspace{1.5cm}para el monitoreo de la fermentaci\'on artesanal del k\'efir de leche \hspace{1.5cm} 
\begin{picture}(0,0) 
\put(20,-7){\includegraphics[width=1.5cm, height=1.5cm]{logos/upiita_logo.png}} \end{picture} 
\underline{\hspace{ \textwidth}}}\hfil\hbox{}}}
\makeatletter
% normal pages
\def\ps@headings{%
\def\@oddhead{\MYhead}%
\def\@evenhead{\MYhead}}%
% title page
\def\ps@IEEEtitlepagestyle{%
\def\@oddhead{\MYhead}%
\def\@evenhead{\MYhead}}%
\makeatother
% make changes take effect
\pagestyle{headings}

% Ajustar espacio entre el encabezado y el cuerpo del texto
\setlength{\headsep}{1.5cm} % Puedes ajustar este valor según tus necesidades



\begin{document}

\begin{titlepage}
    \begin{tikzpicture}[remember picture,overlay]
        % Logo superior izquierdo
        \node[anchor=north west, yshift=-1cm, xshift=1cm] (logoSup)
            at (current page.north west) 
            {\includegraphics[width=3cm]{logos/logo_ipn.png}};

        % Texto al lado del logo superior
        \node[anchor=west, yshift=-0.5cm, xshift=2cm](titulo) at (logoSup.east) 
            {\Large \textbf{Instituto Polit\'ecnico Nacional}};
            
 \draw[line width=1pt] 
  ([xshift=-1.5cm, yshift=-0.5cm] titulo.south west) --
  ([xshift= 1.5cm, yshift=-0.5cm] titulo.south east);

        % Logo inferior izquierdo
        \node[anchor=south west, yshift=3cm, xshift=1cm] 
            at (current page.south west) 
            {\includegraphics[width=3cm]{logos/upiita_logo.png}};

        % Tres líneas negras pegadas al margen izquierdo
        \draw[line width=2pt] (current page.north west) ++(2cm,-5cm) -- ++(0,-18cm);
        \draw[line width=2pt] (current page.north west) ++(2.5cm,-5cm) -- ++(0,-18cm);
        \draw[line width=2pt] (current page.north west) ++(3cm,-5cm) -- ++(0,-18cm);
    \end{tikzpicture}

    % =========================
    % Contenido central de la portada
    % =========================
    \centering

    \vspace*{1.5cm} % bajamos un poco el bloque central

    {\large \textbf{UNIDAD PROFESIONAL INTERDISCIPLINARIA EN INGENIERÍA Y TECNOLOG\'IAS AVANZADAS}}\\[2cm]

    {\Large \textbf{Proyecto Terminal I}}\\[1cm]

    {\huge \textbf{{Prototipo de un sistema telem\'atico para el monitoreo de la fermentaci\'on artesanal del k\'efir de leche}}}\\[1cm]

    {\large Que para obtener el t\'itulo de}\\[0.5cm]
    {\Large \textbf{Ingeniero en Telem\'atica}}\\[1.5cm]

    {\large Presenta:}\\[0.5cm]
    {\large
    Garcia Ledezma Miicxitl Kin}\\[1cm]

    {\large Asesores:}\\[0.5cm]
    {\large
    Dra. Bella Citlali Mart\'inez Seis\\
    Dra. Iclia Villordo Jim\'enez\\
    Ing. Jonas Mart\'inez Lim\'on}\\[1cm]

    {\large \today \par} %esto crea la fecha de hoy
\end{titlepage}

\section*{Resumen}

El presente proyecto tiene como objetivo desarrollar una herramienta tecnológica
que permita a los productores de kéfir artesanal optimizar el tiempo de recolección
del kéfir y determinar si su fermentación se llevó a cabo de manera óptima, con el
fin de garantizar la calidad del producto. Para ello, se desarrollarán cuatro módulos.
El primer módulo se encargará de obtener los datos del proceso de fermentación
mediante sensores. Estos datos serán procesados en el segundo módulo, donde se
realizará un análisis utilizando un modelo basado en reglas de lógica difusa para
evaluar el estado de la fermentación. Luego, el tercer módulo almacenará tanto los
datos obtenidos de los sensores como los resultados del análisis en una base de
datos. Finalmente, el cuarto módulo permitirá visualizar toda esta información a
través de una página web. Este sistema proporcionará a los productores de kéfir
artesanal la posibilidad de analizar sus fermentaciones basándose en datos
estructurados, en lugar de depender exclusivamente de conocimientos empíricos.
\newline
\newline
% Palabras clave
\noindent \textbf{Palabras clave:} kéfir, fermentación, sensores, análisis de datos, lógica difusa, base de datos, página web.

\section*{Abstract}

The objective of this project is to develop a technological tool that allows 
artisanal kefir producers to optimize kefir harvesting time and determine whether 
its fermentation was optimal, in order to ensure product quality. To this end, 
four modules will be developed. The first module will collect data from the fermentation 
process using sensors. This data will be processed in the second module, where an analysis 
will be performed using a fuzzy logic rule-based model to assess the state of the 
fermentation. The third module will then store both the data obtained from the sensors 
and the analysis results in a database. Finally, the fourth module will allow all this 
information to be viewed via a website. This system will provide artisanal kefir producers 
with the ability to analyze their fermentations based on structured data, rather than 
relying solely on empirical knowledge.

% Keywords
\noindent \textbf{Keywords:} kefir, fermentation, sensors, data analysis, fuzzy logic, database, website.
\end{document}
